\documentclass[11pt]{amsart}
\input{../commands.tex}

\newtheorem{hwexercise}{Exercise}

\title{math231br Homework 5 (Due Thursday, 3/5/26)}

\begin{document}
\maketitle

\begin{hwexercise} Discuss: given a vector bundle $E \to X$ over a compact Hausdorff space $X$, can we recover the vector bundle $E$ from the sheaf $\Sec(-,E)$? If so, how? If not, what extra data would you need to rebuild $E$?
\end{hwexercise}


\begin{hwexercise} Prove the fundamental theorem of calculus we used in the de Rham theorem: if $f(x_1, \ldots, x_n)$ is a differentiable function on the unit ball $B \subseteq \R^n$, then
\begin{align*}
    f(x_1, \ldots, x_n) = f(0)+ \sum_{i=1}^n x_i \int_0^1 \frac{df}{dx_i}(tx_1,tx_2, \ldots, tx_n) dt.
\end{align*}
\end{hwexercise}

\begin{hwexercise}
Let $X = \R^2\minus\{(0,0)\}$, and let $H^\ast_\dR(X)$ be the cohomology of the sequence
\begin{align*}
    \R \to \Omega^1_\dR(X) \to \Omega^2_\dR(X) \to 0.
\end{align*}
Prove that $H^1_\dR(X)$ is a 1-dimensional real vector space spanned by
\begin{align*}
    \omega = \frac{x\ dy - y\ dx}{x^2+y^2}.
\end{align*}
\end{hwexercise}

\begin{hwexercise}
Describe the sheaves $\Omega^\ast_\dR$ on the circle $S^1$. What do their sections/stalks look like? What is the cohomology of the sequence
\begin{align*}
    \R \to \Omega^1_\dR(S^1) \to \Omega^2_\dR(S^1) \to \cdots
\end{align*}
and what are generators for the cohomology groups?
\end{hwexercise}


\end{document}
